\section{Preliminaries} \label{sec:preliminaries}

Let $G=\br{V, E}$ be a connected simple graph. Let $\sbr{k}$ be the set of different types of interfaces across the network and let $\lambda\colon V \to 2^{\sbr{k}}$ be a function that assigns to each vertex $v\in V$ its set of available interfaces, such that for each edge $uv \in E$, $\lambda\br{u} \cap \lambda\br{v} \neq \emptyset$.
An assignment of active interfaces is a function $\lambda_A\colon V \to 2^{\sbr{k}}$ such that $\lambda_A\br{v} \subseteq \lambda\br{v}$ for every $v\in V$. We say that an edge $uv \in E$ is \emph{covered} if $\lambda_A\br{u} \cap \lambda_A\br{v} \neq \emptyset$. We denote by $\cov\br{\lambda_A} = \{uv \in E : \lambda_A\br{u} \cap \lambda_A\br{v} \neq \emptyset\}$ the set of edges covered by $\lambda_A$, and by $G_A$ the subgraph of $G$ with the vertex set $\bigcup_{e\in \cov\br{\lambda_A}}e$ and the edge set $\cov\br{\lambda_A}$.

In the \coverage problem, we wish to find an assignment of active interfaces $\lambda_A$ such that $G_A=G$. We call such an assignment a \emph{covering assignment}. In the \connectivity problem, we wish to find an assignment of active interfaces $\lambda_A$ such that, $V\br{G_A}=V$ and $G_A$ is connected. We call such an assignment a \emph{connecting assignment}, i.e., an assignment that induces a connected spanning subgraph of $G$.

We are also given a cost function $c\colon \sbr{k}\times V \to \mathbb{N}_{\geq 0}$, which denotes the cost of activating the interface $i\in\sbr{k}$ at the vertex $v\in V$. As a convention, we assume that $c(i,v)=\infty$ if $i\notin \lambda(v)$. For an assignment $\lambda_A$, define the \emph{cost} of $v$ as:
$$\COST\br{v, \lambda_A} = \sum_{i\in\lambda_A\br{v}}c\br{i, v}.$$
The \emph{max-cost} of an assignment $\lambda_A$ is defined as:
$$\COST \br{\lambda_A}=\max_{v\in V} \COST\br{v, \lambda_A}.$$
The goal of the min-max assignment problem is to find an assignment $\lambda_A$ that minimizes $\COST \br{\lambda_A}$ while ensuring either \coverage or \connectivity. 
The formal problem statements are as follows.

\pbdef{\coverage}{A connected simple graph $G=(V,E)$, available interfaces $\lambda\colon V \to 2^{\sbr{k}}$, a cost function $c\colon \sbr{k}\times V \to \mathbb{N}_{\geq 0}$.}{Find a covering assignment $\lambda_A$ minimizing $\COST\br{\lambda_A}$.}

\pbdef{\connectivity}{A connected simple graph $G=(V,E)$, available interfaces $\lambda\colon V \to 2^{\sbr{k}}$, a cost function $c\colon \sbr{k}\times V \to \mathbb{N}_{\geq 0}$.}{Find a connecting assignment $\lambda_A$ minimizing $\COST\br{\lambda_A}$.}

We denote by $\OPT$ the cost of an optimal solution to the problem at hand. We will refer to the min-max covering assignment problem as \coverage and the min-max connecting assignment problem as \connectivity. Note that there also exists a min-sum variant of these problems, where the goal is to minimize $\sum_{v\in V} \COST\br{v, \lambda_A}$ instead of $\COST\br{\lambda_A}$.
\subsection{Probabilistic tools}

Since our algorithms are randomized, we will need to use the following probability facts in the analysis. 
\begin{theorem}[Markov's Inequality]
Let $X$ be a non-negative random variable and $a > 0$. Then:
\[
    \Pr\sbr{X \geq a} \leq \frac{\mathbb{E}\sbr{X}}{a}.
\]
\end{theorem}

% \begin{theorem}[Wald's Identity \cite{OnCumulativeSumsOfRandomVariables}]
% Let $X_1, X_2, \dots$ be i.i.d.\ random variables with finite expectation $\mathbb{E}[X]$, and let $N$ be a non-negative integer-valued random variable independent of the sequence $(X_i)_{i\geq 1}$ with finite expectation $\mathbb{E}[N]$. Then:
% \[
%     \mathbb{E}\!\left[\sum_{i=1}^{N} X_i\right] = \mathbb{E}[N]\cdot\mathbb{E}[X].
% \]
% \end{theorem}

\begin{theorem}[(Weighted) Chernoff Bound]
    Let $X_1, X_2, \dots, X_n$ be independent random variables taking values in $\sbr{0,1}$, let $w_1, w_2, \dots, w_n$ be non-negative weights, such that for every $i\in \sbr{n}$ we have that $w_i\in \sbr{0,1}$. Let $X=\sum_{i=1}^{n} w_i\cdot X_i$ and let $\mu = \E{X}$. Then for every $\delta_1 > 0$ and $\delta_2 \in (0,1)$ we have that:
    \[
        \Pr\br{X \geq (1+\delta_1)\cdot \mu} \leq \br{\frac{e^{\delta_1}}{(1+\delta_1)^{1+\delta_1}}}^\mu
        \qquad
        \Pr\br{X \leq (1-\delta_2)\cdot \mu} \leq \br{\frac{e^{-\delta_2}}{(1-\delta_2)^{1-\delta_2}}}^\mu.
    \]
\end{theorem}

The proof of the above fact follows the proof of the standard Chernoff bound and we include it in \Cref{sec:appendix} for completeness.

\subsection{Preprocessing}
Before applying our algorithms (for both \coverage and \connectivity), we preprocess the instance in the following way: First, by a polynomial number of guesses and rescaling, we reduce the instances so that all of the costs are rational values in $\sbr{0,1}$ and we have $\OPT\geq 1/2$. Second, we partition vertices into two classes. A vertex $v$ is \emph{cheap} if $\sum_{i\in\lambda\br{v}} c\br{i,v}\leq 1$; for such vertices we may assume all available interfaces are activated. Vertices that are not cheap are \emph{expensive}; for them we may assume $\COST\br{v,\lambda_A}\geq 1$ in the analyzed solutions, while incurring only a factor $2$ loss in the approximation ratio. The full preprocessing procedure and its complete analysis are deferred to the Appendix \ref{sec:appendix-preprocessing}.